\documentclass[10pt, letterpaper]{article}

% Packages:
\usepackage[
    ignoreheadfoot, % set margins without considering header and footer
    top=1.6 cm, % seperation between body and page edge from the top
    bottom=1.6 cm, % seperation between body and page edge from the bottom
    left=2 cm, % seperation between body and page edge from the left
    right=2 cm, % seperation between body and page edge from the right
    footskip=1.0 cm, % seperation between body and footer
    % showframe % for debugging 
]{geometry} % for adjusting page geometry
\usepackage{titlesec} % for customizing section titles
\usepackage{tabularx} % for making tables with fixed width columns
\usepackage{array} % tabularx requires this
\usepackage[dvipsnames]{xcolor} % for coloring text
\definecolor{primaryColor}{RGB}{0, 79, 144} % define primary color
\usepackage{enumitem} % for customizing lists
\usepackage{fontawesome5} % for using icons
\usepackage{amsmath} % for math
\usepackage[
    pdftitle={Leo Yao's Resumé},
    pdfauthor={Leo Yao},
    pdfcreator={LaTeX},
    colorlinks=true,
    urlcolor=primaryColor
]{hyperref} % for links, metadata and bookmarks
\usepackage[pscoord]{eso-pic} % for floating text on the page
\usepackage{calc} % for calculating lengths
\usepackage{bookmark} % for bookmarks
\usepackage{lastpage} % for getting the total number of pages
\usepackage{changepage} % for one column entries (adjustwidth environment)
\usepackage{paracol} % for two and three column entries
\usepackage{ifthen} % for conditional statements
\usepackage{needspace} % for avoiding page brake right after the section title
\usepackage{iftex} % check if engine is pdflatex, xetex or luatex

% Ensure that generate pdf is machine readable/ATS parsable:
\ifPDFTeX
    \input{glyphtounicode}
    \pdfgentounicode=1
    % \usepackage[T1]{fontenc} % this breaks sb2nov
    \usepackage[utf8]{inputenc}
    \usepackage{lmodern}
\fi



% Some settings:
\AtBeginEnvironment{adjustwidth}{\partopsep0pt} % remove space before adjustwidth environment
\pagestyle{empty} % no header or footer
\setcounter{secnumdepth}{0} % no section numbering
\setlength{\parindent}{0pt} % no indentation
\setlength{\topskip}{0pt} % no top skip
\setlength{\columnsep}{0cm} % set column seperation
\makeatletter
\let\ps@customFooterStyle\ps@plain % Copy the plain style to customFooterStyle
\patchcmd{\ps@customFooterStyle}{\thepage}{
    \color{gray}\textit{\small Leo Yao - Page \thepage{} of \pageref*{LastPage}}
}{}{} % replace number by desired string
\makeatother
\pagestyle{empty}

\titleformat{\section}{\needspace{2\baselineskip}\bfseries\large}{}{0pt}{}[\vspace{1pt}\titlerule]

\titlespacing{\section}{
    % left space:
    -1pt
}{
    % top space:
    0.18 cm
}{
    % bottom space:
    0.12 cm
} % section title spacing

\renewcommand\labelitemi{$\circ$} % custom bullet points
\newenvironment{highlights}{
    \begin{itemize}[
        topsep=0.06 cm,
        parsep=0.06 cm,
        partopsep=0pt,
        itemsep=0pt,
        leftmargin=0.4 cm + 10pt
    ]
}{
    \end{itemize}
} % new environment for highlights

\newenvironment{highlightsforbulletentries}{
    \begin{itemize}[
        topsep=0.10 cm,
        parsep=0.10 cm,
        partopsep=0pt,
        itemsep=0pt,
        leftmargin=10pt
    ]
}{
    \end{itemize}
} % new environment for highlights for bullet entries


\newenvironment{onecolentry}{
    \begin{adjustwidth}{
        0.2 cm + 0.00001 cm
    }{
        0.2 cm + 0.00001 cm
    }
}{
    \end{adjustwidth}
} % new environment for one column entries

\newenvironment{twocolentry}[2][]{
    \onecolentry
    \def\secondColumn{#2}
    \setcolumnwidth{\fill, 5 cm}
    \begin{paracol}{2}
}{
    \switchcolumn \raggedleft \secondColumn
    \end{paracol}
    \endonecolentry
} % new environment for two column entries

\newenvironment{header}{
    \setlength{\topsep}{0pt}\par\kern\topsep\centering\linespread{1.5}
}{
    \par\kern\topsep
} % new environment for the header

% \newcommand{\placelastupdatedtext}{% \placetextbox{<horizontal pos>}{<vertical pos>}{<stuff>}
%   \AddToShipoutPictureFG*{% Add <stuff> to current page foreground
%     \put(
%         \LenToUnit{\paperwidth-2 cm-0.2 cm+0.05cm},
%         \LenToUnit{\paperheight-1.0 cm}
%     ){\vtop{{\null}\makebox[0pt][c]{
%         \small\color{gray}\textit{Last updated in September 2024}\hspace{\widthof{Last updated in September 2024}}
%     }}}%
%   }%
% }%

% save the original href command in a new command:
\let\hrefWithoutArrow\href

% new command for external links:
\renewcommand{\href}[2]{\hrefWithoutArrow{#1}{\ifthenelse{\equal{#2}{}}{ }{#2 }\raisebox{.15ex}{\footnotesize \faExternalLink*}}}


\begin{document}
    \newcommand{\AND}{\unskip
        \cleaders\copy\ANDbox\hskip\wd\ANDbox
        \ignorespaces
    }
    \newsavebox\ANDbox
    \sbox\ANDbox{}

    % \placelastupdatedtext
    \begin{header}
        \textbf{\fontsize{24 pt}{24 pt}\selectfont Leo Yao}

        \vspace{0.3 cm}

        \normalsize
        \mbox{{\color{black}\footnotesize\faMapMarker*}\hspace*{0.13cm}Palo Alto, CA}%
        \kern 0.25 cm%
        \AND%
        \kern 0.25 cm%
        \mbox{\hrefWithoutArrow{mailto:leoyao@andrew.cmu.edu}{\color{black}{\footnotesize\faEnvelope[regular]}\hspace*{0.13cm}leoyao@andrew.cmu.edu}}%
        \kern 0.25 cm%
        \AND%
        \kern 0.25 cm%
        \mbox{\hrefWithoutArrow{tel:+1650-285-0101}{\color{black}{\footnotesize\faPhone*}\hspace*{0.13cm}+1 (650) 285-0101}}%
        \kern 0.25 cm%
        \AND%
        % \kern 0.25 cm%
        % \mbox{\hrefWithoutArrow{https://yourwebsite.com/}{\color{black}{\footnotesize\faLink}\hspace*{0.13cm}yourwebsite.com}}%
        % \kern 0.25 cm%
        % \AND%
        \kern 0.25 cm%
        \mbox{\hrefWithoutArrow{https://linkedin.com/in/leowo/}{\color{black}{\footnotesize\faLinkedinIn}\hspace*{0.13cm}leowo}}%
        \kern 0.25 cm%
        \AND%
        \kern 0.25 cm%
        \mbox{\hrefWithoutArrow{https://github.com/LeoY20}{\color{black}{\footnotesize\faGithub}\hspace*{0.13cm}LeoY20}}%
    \end{header}

    \vspace{0.3 cm - 0.3 cm}

    % \section{Welcome to RenderCV!}

    %     \begin{onecolentry}
    %         \href{https://rendercv.com}{RenderCV} is a LaTeX-based CV/resume version-control and maintenance app. It allows you to create a high-quality CV or resume as a PDF file from a YAML file, with \textbf{Markdown syntax support} and \textbf{complete control over the LaTeX code}.
    %     \end{onecolentry}

    %     \vspace{0.2 cm}

    %     \begin{onecolentry}
    %         The boilerplate content was inspired by \href{https://github.com/dnl-blkv/mcdowell-cv}{Gayle McDowell}.
    %     \end{onecolentry}
    
    % \section{Quick Guide}

    % \begin{onecolentry}
    %     \begin{highlightsforbulletentries}


    %     \item Each section title is arbitrary and each section contains a list of entries.

    %     \item There are 7 unique entry types: \textit{BulletEntry}, \textit{TextEntry}, \textit{EducationEntry}, \textit{ExperienceEntry}, \textit{NormalEntry}, \textit{PublicationEntry}, and \textit{OneLineEntry}.

    %     \item Select a section title, pick an entry type, and start writing your section!

    %     \item \href{https://docs.rendercv.com/user_guide/}{Here}, you can find a comprehensive user guide for RenderCV.

    %     \end{highlightsforbulletentries}
    % \end{onecolentry}

    \section{Education}
        \begin{twocolentry}{
        % two-col takes in 1 arg for the right one and dynamic left.
        \text{August 2024 – May 2028}}
            \textbf{Carnegie Mellon University}
        \end{twocolentry}
        % degree, dates
        \begin{twocolentry}{
            \text{Pittsburgh, PA}}
            \textit{B.S. in Computer Science and Statistics \& Machine Learning}
        \end{twocolentry}
        
        \vspace{0.06 cm}
        \begin{onecolentry}
            \begin{highlights}
                \item \textbf{QPA:} 3.65/4.00 
                \item \textbf{Relevant Coursework:} Computer Systems, Parallel \& Sequential Data Structures \& Algorithms, Computer Security, Generative AI (Graduate), Machine Learning, Statistical Computing, Functional Programming, Theoretical Computer Science
                \item \textbf{Awards:} Dean's List, High Honors
            \end{highlights}
        \end{onecolentry}

    \section{Experience}
        \begin{twocolentry}{\text{May 2026 - August 2026}}
            \textbf{Software Developer Intern}
        \end{twocolentry}
        \begin{twocolentry}{\text{San Jose, CA}}
            \textit{IBM}
        \end{twocolentry}
        \begin{onecolentry}
            \begin{highlights}
                \item Built an AI-agent observability connector for watsonx.governance, integrating three vendor APIs (Arrow Flight/gRPC, GraphQL, REST) to normalize per-agent cost, latency, error-rate, and LLM evaluator scores.
                \item Instrumented multi-agent LangGraph systems on two observability platforms (Arize AX, Phoenix) with OpenTelemetry/OpenInference, emitting all 8 span kinds plus PII/guardrail spans.
                \item Cut peak worker memory 6.4$\times$ by chunking unbounded, customer-controlled span exports, and reduced live-path API calls 22$\rightarrow$1 ($\sim$95\%) by redesigning around a $\sim$1,000 query-complexity cap.
                \item Wrote 373 tests at $\sim$2:1 test-to-source, including a differential-oracle suite that recomputes every metric independently over 1,500 randomized inputs and 16,000 percentile comparisons.
            \end{highlights}
        \end{onecolentry}
        \begin{twocolentry}{\text{January 2026 - May 2026}}
            \textbf{Teaching Assistant}
        \end{twocolentry}
        \begin{twocolentry}{\text{Pittsburgh, PA}}
            \textit{Carnegie Mellon University}
        \end{twocolentry}
        \vspace{0.06 cm}
        \begin{onecolentry}
            \begin{highlights}
                \item Led weekly 1.5-hour recitations and 3 hours of office hours for 20+ students in 15-150 Principles of Functional Programming, CMU's required course on SML, recursion, and proofs of program correctness.
                \item Graded 300+ homework and exam submissions weekly in SML, giving technical feedback on correctness, efficiency, and style across recursive, higher-order, and parallel implementations.
            \end{highlights}
        \end{onecolentry}

        \begin{twocolentry}{
            \text{September 2025 - May 2026}}
            \textbf{Research Assistant}
        \end{twocolentry}
        \begin{twocolentry}{
            \text{Pittsburgh, PA}}
            \textit{Carnegie Mellon University}
        \end{twocolentry}
        \vspace{0.06 cm}
        \begin{onecolentry}
            \begin{highlights}
                \item Owned the training and evaluation harness for quantization-aware, FPGA-bound GNNs at CERN CMS: resumable checkpointing, CLI-driven hyperparameter sweeps, and evaluation against the PUPPI baseline.
                \item Adapted lightweight PyTorch transformer models for use in particle clustering on CERN HGCAL data, creating a custom data pipeline with NumPy and Awkward to properly load in data for training.
            \end{highlights}
        \end{onecolentry}
        % \begin{twocolentry}{
        % \text{June 2025 – August 2025}}
        %     \textbf{Research Assistant, SPICE Lab}
        % \end{twocolentry} 
        % \begin{twocolentry}{
        %     \text{Pittsburgh, PA}}
        % \textit{Carnegie Mellon University}
        % \end{twocolentry}
        % \vspace{0.10 cm}
        % \begin{onecolentry}
        %     \begin{highlights}
        %         \item Engineered a Python data pipeline (Pandas, NumPy, scikit-learn) to clean, normalize, and extract features from structured energy datasets.
        %         \item Optimized PyTorch neural networks using LightGBM feature selection and hyperparameter tuning, improving accuracy by 75\% over initial baseline.
        %     \end{highlights}
        % \end{onecolentry}

        % \vspace{0.2 cm}
        
        % \begin{twocolentry}{
        % \text{July 2023 - August 2023}}
        %     \textbf{Data Science Researcher}
        % \end{twocolentry}
        % \begin{twocolentry}{
        %     \textit{Irvine, CA}}
        %     \textit{UC Irvine COSMOS Program}
        % \end{twocolentry}

        % \vspace{0.10 cm}
        % \begin{onecolentry}
        %     \begin{highlights}
        %         \item Implemented and benchmarked several regression models (Linear, LASSO, Random Forest, etc.) to evaluate accuracy and performance in predicting brain grey matter from measured human physical features.
        %         \item Plotted results from each model with ggplot2, summarizing results in a project poster; presented poster at a research symposium to peers in the program and UCI faculty.  Poster and code available upon request.
        %     \end{highlights}
        % \end{onecolentry}
    \section{Projects}
        % shortened version
        \begin{twocolentry}{\text{February 2026}}
            \textbf{PharmaSentinel} $|$ \textit{Python, PostgreSQL, Supabase, MCP}
        \end{twocolentry}
        \vspace{0.06 cm}
        \begin{onecolentry} 
            \begin{highlights}
                \item Engineered a multi-agent backend system in Python, orchestrating LLM and openFDA REST APIs to automate real-time supply chain monitoring and substitute identification.
                \item Architected event-driven PostgreSQL databases utilizing Supabase Realtime to detect inventory state changes, automatically triggering agent pipelines and routing resolutions to downstream clients.
            \end{highlights}            
        \end{onecolentry}
        \begin{twocolentry}{
        \text{May 2025 – August 2025}}
            \textbf{Computer Systems} $|$ \textit{C, POSIX threads}
        \end{twocolentry}

        \vspace{0.06 cm}
        \begin{onecolentry}
            \begin{highlights}
                % \item \textbf{Cache Simulator} Developed a write-back/write-allocate cache simulator with LRU eviction policy that parses instruction traces, accurately modeling hit/miss/eviction rates across various cache geometries.
                \item \textbf{Dynamic Memory Allocator}: Developed a dynamic memory allocator in C for Linux. Achieved 74.3\% memory utilization and a throughput of 15885 kilo-operations per second, ranking at the top of the class.
                \item \textbf{Tiny Linux Shell}: Engineered a Linux shell featuring foreground/background job control, I/O redirection, and robust signal handling to manage concurrent process execution.
                % \item \textbf{Web Proxy Server}: Created a multithreaded web proxy server that facilitates communication between clients and web servers.
                \item \textbf{Shark File System}: Built a thread-safe concurrent file system with file renaming and descriptor position management, using mutex synchronization for multithreaded correctness.
            \end{highlights}
        \end{onecolentry}
        % \begin{twocolentry}{\text{July 2025}}
        %     \textbf{Shark File System}
        % \end{twocolentry}
        % \vspace{0.10 cm}
        % \begin{onecolentry}
        %     \begin{highlights}
        %         \item Developed a thread-safe concurrent FAT-style file system in C with file creation, deletion, renaming, reading, and writing capabilities. Optimized performance/correctness 
        %         with mutexes to mitigate race conditions.
        %     \end{highlights}
        % \end{onecolentry}

    % \section{Leadership}
    %     \begin{twocolentry}{
    %         \text{November 2024 – Present}}
    %         \textbf{Joint Funding Committee Member}
    %     \end{twocolentry}
    %     % \begin{twocolentry}{
    %     %     \textit{Pittsburgh, PA}}
    %     %     \textit{CMU Student Government}
    %     % \end{twocolentry}

    %     \vspace{0.10 cm}
    %     \begin{onecolentry}
    %         \begin{highlights}
    %             \item Managing the distribution and allocation of approximately \$2.1 million to 300+ student organizations.
    %         \end{highlights}
    %     \end{onecolentry}    
    % \section{Awards}
    %     \begin{onecolentry}
    %         \textbf{Dean's List, High Honors} (Spring 2025)
    %     \end{onecolentry}
    \section{Skills}
        \begin{onecolentry}
            \textbf{Programming Languages:} Python, C, C++, SQL, SML, x86-64 Assembly \newline
            \textbf{Backend \& APIs:} FastAPI, asyncio, Redis, GraphQL, REST, PostgreSQL, Supabase, Docker, Linux \newline
            \textbf{AI \& Observability:} LangGraph, OpenTelemetry, OpenInference, Arize AX/Phoenix, MCP \newline
            \textbf{Data \& Testing:} pandas, NumPy, PyTorch, Apache Arrow Flight, pytest, Git, GDB, Valgrind
        \end{onecolentry}
    \vspace{0.10 cm}

    

    

\end{document}